\paragraph{The integral lower bound on a single component.}
\hypertarget{source-slice-metrics}{}
Keep the integer \(h\) of \eqref{eq:chow-lower-section}. A product
of \(h\) integral light forms, and any integral sum of those
products, satisfies the integral framed-point test at weight
\(h w_L\). This follows by multiplying their values in each framed
valuation ring. Positive weight and cuspidality persist under this
multiplication. Suppose the resulting form \(g\) restricts
nontrivially to \(Z\). On the generic fibre multiply it by the
idempotent of \(Z\). The resulting section \(g_Z\) is algebraic
and cuspidal. In every framed DVR test its generic point belongs
to exactly one component. The value tested there equals the
original value or zero. Thus \(g_Z\) belongs to the same global
integral lattice and is nonzero. This proof also covers components
whose closures meet in a special fibre.

\begin{proposition}[The lower norm bound on a single component]\label{prop:integral-chow-component-height}
For the fixed exponent \(h\) of \eqref{eq:chow-lower-section}
let \(g\) be an integral sum of products of \(h\) integral
light sections with \(g|_Z\ne0\).
Use the constant \(P_0\) of \eqref{eq:component-chow-lower-constant},
chosen from the common column-size and weight data before \(g\) and \(Z\).
On every such component,
\[
 \operatorname{mean}_\sigma\log\|g|_Z\|_{\sup,\sigma}\ge-hP_0S_L.
\]

The integral lifting uses the generic component idempotent
and the original cusp lattice at weight \(h w_L\).
\end{proposition}

\begin{proof}
\proofstep{Bounding the supremum by algebraic evaluations.}
At an embedding \(\sigma\), the supremum norm of \(g_Z\) is bounded
by the supremum of the norms of \(g\), with the same metric, under embeddings \(\Omega_Z\to\mathbb C\) extending
\(\sigma\). Indeed, work on an open affine flag chart over \(K\)
on which the frame has been trivialized. Its ring is countable.
For a nonzero element of the ring the set of complex points
where it vanishes has empty interior in each analytic component.
Each geometric component dominates the integral chart over \(K\),
and on the smooth complex chart the assertion follows from
the identity theorem. Baire's theorem therefore gives a dense
set of points where the ring embeds in \(\mathbb C\); the map
at such a point extends to its fraction field.

Equivalently the evaluation map on the chart ring \(A\) is injective at each point of the dense set
\[
  \bigcap_{a\in A\setminus\{0\}}\bigl\{x:a(x)\ne0\bigr\}.
\]
Let \(F\) be the coefficient of the section in an adapted frame and
let \(J>0\) be the norm of that frame. The section has norm \(|F|J\).
Under a change of frame with weight character \(\chi\), the coefficient
and frame norm transform inversely, so
\[
  |\chi|\,|F|\cdot |\chi|^{-1}J=|F|J.
\]
Consequently the section norms on the dense set agree
with the generic norm. The metric is continuous across the
chart and the corresponding flag cells are dense in the component.
At conjugate embeddings one uses the conjugate uniformization
with its prescribed metric. This proves the supremum comparison.

Apply \eqref{eq:slice-integral-norm} for nonzero integral
\(F\) of the required weight \(v\). In the notation of that
estimate, put \(P_*=3000Z^2\) with \(Z=fr+L+\log d_E+|\Omega|\).
Then
\[
  \operatorname{mean}_\sigma\log\|F\|_{\sup,\sigma}
        \ge-P_* (1+\log|v|)|v| .
\]
Here \(|v|\) is the weight size. Apply the bound to \(g_Z\)
with \(v=h w_L\) and divide by \(hS_L\). Use \(1\le |h w_L|\le hS_L\)
and
\[
 \log h\le \log\!\big(8(s+2)^3(2+P_O)\big)
            +2\log N+\log\delta .
\]
The family bounds are \(\log N\le n_B\) and
\(\log\delta\le1041z^3\), where \(n_B\) bounds the logarithm
of the column size. One may take
\begin{equation}\label{eq:component-chow-lower-constant}
 P_0=\max\{1,P_*[1+\log S_L+\log(8(s+2)^3(2+P_O))
                                         +2n_B+1041z^3]\}.
\end{equation}
Thus the lower bound of Proposition~\ref{prop:slice-lower-integral-norm}
at the single weight \(h w_L\) gives \eqref{eq:chow-lower-section}.
The integer \(h\) and the bound for \(P_0\) are common to all components. \label{cor:slice-metrics-1}

\end{proof}

\paragraph{From algebraic integral points to the orders of the slice.}
\label{par:slice-valued-orders}
The depth test available on algebraic points is a bound on the minimum valuation of the light sections in an adapted integral frame. We explain how
it gives orders at the nonalgebraic places used by the height.
The following extension lemma is the input to the branch calculation.

\begin{lemma}[Integral tests at the closed fibre]\label{lem:slice-valuative-extension}
Let \(O\) be a complete
DVR of a number field completion, with uniformizer \(\pi\).
Let \(A\) be a flat normal integral algebra of finite type over
\(O\), and let \(g\in\operatorname{Frac}A\). Suppose that in
every map \(A\to V\) to a finite extension DVR of \(O\) whose
generic point belongs to a fixed dense open where \(g\) is
defined, the value of \(g\) lies in \(V\). Then the simultaneous local conclusion is
\[
 g\in\bigcap_{\substack{q\in\operatorname{Spec}A\\ q+(\pi)\ne A}}A_q.
\]


\end{lemma}

\begin{proof}
A pole of \(g\) at a height-one prime \(P\subset q\)
gives a point \((P,\infty)\) on the closure of its graph
in \(\operatorname{Spec} A\times\mathbb P^1\), by evaluating
on the DVR at \(P\). Taking the closure gives a closed point
\((y,\infty)\) on the special fibre, since 
\[
  (q,\pi)\ne A.
\]
An integral curve through that point with generic point in
the prescribed open can be chosen as follows.

In the local
domain of the graph at this point, of dimension \(d\), choose
a nonzero \(b\) divisible by \(\pi\) such that \(D(b)\) is in
the original graph over that open and \(1/g\) is nonzero
there. The case \(g=0\) is immediate. Complete \(b\) to a
system of parameters \(b,a_2,\ldots,a_d\). Write \(B\) for this local domain and \(I=(a_2,\ldots,a_d)\).
Killing one element of the maximal ideal lowers the dimension
by at most one \cite{slice-dimension-drop}. Hence
\[
 \dim B/I\ge d-(d-1)=1,\qquad \sqrt{I+(b)}=\mathfrak m_B .
\]
If every minimal prime over \(I\) contained \(b\), these two facts would
give \(\dim B/I=0\). Choose such a prime \(\mathfrak p\) avoiding \(b\).
In the local domain \(B/\mathfrak p\) the nonzero image of \(b\) generates a
maximal-primary ideal. The same dimension inequality now gives
\(\dim B/\mathfrak p=1\).

Let \(C\) be the closure of this branch. It dominates \(\operatorname{Spec} O\)
since \(\pi\mid b\). At the chosen closed point its local dimension is one
and its residue field is finite over the residue field of \(O\).
The dimension formula over the universally catenary DVR \(O\) gives
\[
 1=1+\operatorname{trdeg}_{\operatorname{Frac}O}K(C)
       -\operatorname{trdeg}_{O/(\pi)}\kappa(y).
\]
Thus the finitely generated field \(K(C)\) is finite over
\(\operatorname{Frac} O\) \cite{dim-slice}. A local ring of the
normalization above the closed point is a DVR dominating \(O\).
The extension of the valuation of a complete DVR to this finite
extension is unique and its valuation ring is finite over \(O\)
\cite{dvr-slice}. On the resulting test the generic point lies in
\(D(b)\), whereas \(1/g\) has strictly positive valuation.
This contradicts the assumption on \(g\).
Consequently \(g\) has no pole at any height-one prime of \(A_q\).
Normality now gives
\[
 g\in\bigcap_{\operatorname{ht}_{A_q}\mathfrak p=1}
               (A_q)_{\mathfrak p}=A_q
\]
by the intersection theorem for Noetherian normal domains
\cite{slice-normal-intersection}.\label{cor:slice-metrics-2}

\end{proof}

\begin{proposition}[Transfer of the framed depth and metric tests]\label{prop:slice-chart-metric-readings}
Use the selected light columns and, before the final lift,
the auxiliary \(\eta\) with its stated pointwise depth
inequality. On the chosen normalized charts the inequality
holds at the generic slice branches, using the coefficient
baseline on vertical branches and without that term on
horizontal branches. At a horizontal branch \(q\) its incidence exponent obeys
\[
 0\le E_q\le\frac{c\,n_q(\eta)}{2i_\eta}.
\]

Evaluation in the adapted complex frames gives the paired light and
auxiliary metric bound in the proof. At the definite lift
only the finite baseline is used.
\end{proposition}

\begin{proof}
\proofstep{Transferring the valuation test to the flag chart.}
Apply this assertion on integral framed charts which record a
smallest light column \(w_b\). Such a chart adjoins the regular
ratios \(w_l/w_b\) and uses the normalization of the corresponding
blowup chart. For a valued point with valuation \(v\), choose a column with least
finite order. Then
\[
 v(w_l/w_b)=v(w_l)-\min_j v(w_j)\ge0
\]
for every \(l\), so all the adjoined ratios are in this valuation ring. The required cover uses the
integral model, with frames chosen over its affine opens. To see that these charts suffice for the centre criterion, first
take a vertical DVR. Complete it. Its map from the coefficient
DVR extends to the completion, and its order on the original
field stays unchanged. Use the proper model at this prime
and its flag bundle as in the construction of \(y_\zeta\).

At each elimination step in the flag matrix, choose a nonzero eligible pivot \(a\) of minimum valuation among the eligible entries \(a_j\). Comparability of principal ideals in the valuation ring \(V\) gives
\[
        v(a_j)\ge v(a)\quad\Longrightarrow\quad a_j/a\in V .
\]
The elimination coordinates therefore give an integral point on this flag chart.
The theta generator and these flags may be chosen on finitely
many trivializing charts. After adjoining the ratios for a
smallest light column, the generic fibre over the original
nonboundary locus is still covered: the columns there
generate the unit ideal in a trivialization.

Inverting the
pivot therefore gives that chart on the generic fibre.
After passing to the completed coefficient ring take the closures
of the generic components and normalize. These normalizations
are finite, since the base is excellent, and the valuative
criterion for the resulting finite maps lifts the point with
its generic lift. On the open generic flag charts the
normalization agrees with the smooth model. The resulting
affines are normal and flat over the DVR by torsion-freeness.
In particular if \(A\) is such an affine algebra over the coefficient DVR, its flatness follows from
\[
  0\longrightarrow A\xrightarrow{\ \pi\ }A\longrightarrow A/\pi A\longrightarrow0.
\]

\proofstep{Horizontal orders and the line-frame correction.}
For a horizontal valuation there are two orders to keep track
of. Let \(T\) be its DVR with a copy of \(K\) on which its
valuation is trivial. One can extend it, with ramification
index one, to a DVR \(V\) containing \(K_\lambda\) as a
field of units together with zero. To see this, complete \(T\)
and choose a coefficient field containing \(K\). Explicitly,
lift a transcendence basis of the residue field \(k/K\).
The embedding of the rational function field so obtained
extends across the separable algebraic residue extension by
Hensel lifting, uniquely with the chosen residues. Thus

\[
  \widehat T=k[[u]]
\]
 for a uniformizer \(u\).
Embed \(k\) and \(K_\lambda\) over \(K\) in a common field \(k'\), using the fraction field of a prime quotient of \(k\otimes_K K_\lambda\). The resulting inclusions are
\[
       T\longrightarrow\widehat T=k[[u]]
          \longrightarrow V=k'[[u]],\qquad K_\lambda\subset k'.
\]
The first nonzero coefficient of a series stays nonzero after extending \(k\) to \(k'\), so \(u\)-adic orders are unchanged. Every nonzero element of \(K_\lambda\) is a constant unit of \(V\).

Let \(R'\subset k'\) be a valuation ring extending the ring of \(K_\lambda\).
The composite valuation has ring and maximal ideal
\[
 W=\{a\in V:\bar a\in R'\},\qquad
 \mathfrak m_W=\{a\in V:\bar a\in\mathfrak m_{R'}\}.
\]
Thus \(\mathfrak m_V\subset\mathfrak m_W\). The
coefficient uniformizer \(\pi\) is a unit of \(V\), but its
residue in \(k'\) lies in \(\mathfrak m_{R'}\), so it belongs
to \(\mathfrak m_W\).
In particular the coefficient uniformizer distinguishes the two valuation rings by
\[
  \pi\in V^\times,\qquad \pi\in\mathfrak m_W .
\]
The extension of valuations used here follows by domination of
a local ring by a valuation ring, or equivalently the extension
theorem for Krull valuations.

Apply proper extension, flag pivots, light pivots and the normal
charts first over \(W\). The morphism to the fine valuation ring
evaluates the chart in \(W\subset V\). Pullback to \(V\) gives
the original horizontal point after the unramified extension.
The integral frames differ there by units on the diagonals:
the two underlying flag points agree by separatedness, and both
frames are bases over \(V\). In this chart the inverse image of
\(\mathfrak m_W\) contains both the inverse image of
\(\mathfrak m_V\) and \(\pi\).
Writing \(A\) for the chart ring, the two inverse images thus satisfy
\[
  (\mathfrak m_V\cap A)+\pi A\subseteq\mathfrak m_W\cap A .
\]
 This proves the specialization condition for the chosen model. An arbitrary neighborhood of a horizontal centre need not satisfy it.
The good-level compactification used here is the same at each
run and at all its places. \label{cor:slice-metrics-3}

\proofstep{The corresponding least orders of the selected sections.}
The algebraic test \eqref{eq:slice-tied-depth}
in the light and auxiliary frames gives, with
\(v(\pi_\lambda)=1\) for a uniformizer of \(K_\lambda\),
\[
 0\le v(\eta),\qquad
 v(w_b)\le\nu_\lambda+c'v(\eta),\qquad c'=c/(2i_\eta).
\]
Take \(v(\eta)\) finite by excluding its zero divisor in the
open used to test regularity. It is nonzero on each component.
For each integer \(l\ge1\) apply the assertion to
\[
 g_l=\pi^{\lceil e_O l\nu_\lambda\rceil}
           \eta^{\lceil l c'\rceil}w_b^{-l},
 \qquad e_O=\operatorname{ord}_O(\pi_\lambda).
\]
Use the readings in the chosen trivializations. The algebraic
test and \(v(\pi)=1/e_O\) make every such value integral.
Apply the same criterion to \(\eta\). Evaluate the resulting
local-ring expressions at the centre of a slice valuation,
where the generic readings are defined. After evaluation at the branch, the inequalities are
\[
 l\,v(w_b)\le
   \lceil e_Ol\nu_\lambda\rceil\,v(\pi)
      +\lceil lc'\rceil\,v(\eta).
\]
Divide by \(l\), keeping the branch fixed. The rounding
errors tend to zero. On a vertical branch
\(e_Ov(\pi)=1\) in the \(v(\pi_\lambda)=1\) normalization,
which retains the coefficient baseline.

On the horizontal
branch the coefficient valuation is trivial and
\(v(\pi)=0\), so only \(c'v(\eta)\) remains. Multiplying
by the ramification factor on a vertical branch converts to
its integer order. Summing with residue degrees gives exactly
the auxiliary coefficient in the finite and horizontal
comparisons for \(H_L\). Indeed 
\[
  E_q\le c'n_q(\eta)
\]
and each \(n_q(\eta)\ge0\). For a horizontal polynomial
factor \(q\) its sphere height may be negative. Applying
\eqref{eq:slice-sphere-lower} to the nonnegative difference
\((c/i_\eta)n_q(\eta)-E_q\) incurs altogether at most
\(s(c/i_\eta)a_\eta\ell_N\). This is the term
\(s\beta a_e\ell_N\) above. At the last lift the
test has only its baseline and the power of \(\eta\) is omitted.

\proofstep{The complex metric in the adapted frame.}
For the complex metric, choose the adapted frame attached to a
generic embedding of \(\Omega_Z\). The invariant norms of the
light forms then agree with the chart norms. For the auxiliary
section the generic norm is at most the chart norm. The chart
estimates with first-minimum function \(P\) are
\[
 -\log\big(\sum_l\|w_l\|^2\big)^{1/2}\le P_5 D_F+cP,
 \qquad \log\|\eta\|\le-i_\eta(P-P_2).
\]
Since \(c\ge0\) and \(i_\eta>0\), substitution gives at each
generic reading
\[
 -\log\big(\sum_l\|w_l\|^2\big)^{1/2}
  \le P_5D_F+ (c/i_\eta)(-\log\|\eta\|)+cP_2 .
\]
The auxiliary reading is nonzero under a field embedding. These
arguments reduce the metric comparisons to the algebraic-point
bound and the two analytic first-minimum
bounds on the selected forms. \label{cor:slice-metrics-4}
\end{proof}
